\scanpage{115}
On trouve l'identité \deleted{\unsure{des|dès qu'}} on trouve comme \unsure{isom|isom.} $g_M$
\[
X_M\times X_M\longrightarrow\widehat{\Sigma}_M
\]
celui qui, pour $\alpha,\beta\in\Gamma\unsure{\gamma|\mathfrak g}M$, définissant des sections de $X$ (i.e. tels que $\alpha,\beta$ aient tr et det nuls, et sont $\ne0$ sur tte fibre) distinctes en chaque fibre (i.e.
\[
\operatorname{tr}(\alpha\beta)\ne0\qquad\text{sur chaque fibre),}
\]
associe la \deleted{élément}\inserted{section} de $\Sigma\subset\inserted{\gamma M}$ définie par
\[
-[\alpha,\beta]/\unsure{\operatorname{tr}|\varphi}(\alpha\unsure{,}\beta)
\]
(et non $+[\alpha,\beta]/\operatorname{tr}\alpha\beta$). Il faut vérifier sur mes propres textes \unsure{canoniques|cosmiques|conséquente} \ldots. Si on veut l'interpréter comme section de $\widehat\Sigma$, il faut prendre $-[\alpha,\beta]/\operatorname{tr}(\alpha\beta)+1$ ou encore $\operatorname{tr}(\alpha\beta)\cdot1-[\alpha,\beta]$. Donc l'iso \deleted{$\unsure{X_M\times X}$}
\[
X_M\times X_M\unsure{-|\longrightarrow}\widehat{\Sigma}_M
\]
\deleted{semble} provient d'un \deleted{iso} \inserted{hom} \deleted{\unsure{sur}} au niveau des cônes \deleted{\unsure{épointés|associés}} affines $C=\widetilde V_{\mathrm{aff}}(q_{\gamma M})$ et $D=V_{\mathrm{aff}}(q_M)$ \deleted{$\widetilde X_M\times\widetilde X_M\longrightarrow\unsure{\widetilde\Sigma}$}
% archive sheet number: 114
